%%%%%%%%%%%%%%%%%%%%%%%%%%%%%%%%%%%%%%%%%%%%%%%%%%%%%%%%%%%%
% (01) Input Settings
%%%%%%%%%%%%%%%%%%%%%%%%%%%%%%%%%%%%%%%%%%%%%%%%%%%%%%%%%%%%
\input{settings.tex}

%%%%%%%%%%%%%%%%%%%%%%%%%%%%%%%%%%%%%%%%%%%%%%%%%%%%%%%%%%%%
% (02) DOCUMENT INFORMATION - FROM YAML METADATA
%%%%%%%%%%%%%%%%%%%%%%%%%%%%%%%%%%%%%%%%%%%%%%%%%%%%%%%%%%%%
\newcommand{\DocTitle}{گرافیک}

\newcommand{\CourseTitle}{برنامه‌سازی پیشرفته و کارگاه}

\newcommand{\Semester}{بهار ۱۴۰۵}

\newcommand{\FirstProfessor}{دکتر روح‌الله احمدیان}

\newcommand{\SecondProfessor}{}

\newcommand{\Writers}{سام قربانی، نسترن افسری، نوا جغتائی}

\newcommand{\Faculty}{دانشکده ریاضی و علوم کامپیوتر}

\newcommand{\University}{دانشگاه صنعتی امیرکبیر}

%%%%%%%%%%%%%%%%%%%%%%%%%%%%%%%%%%%%%%%%%%%%%%%%%%%%%%%%%%%%
% DOCUMENT BODY
%%%%%%%%%%%%%%%%%%%%%%%%%%%%%%%%%%%%%%%%%%%%%%%%%%%%%%%%%%%%
\begin{document}
	
%%%%%%%%%%%%%%%%%%%%%%%%%%%%%%%%%%%%%%%%%%%%%%%%%%%%%%%%%%%%
% === TITLE PAGE ===
%%%%%%%%%%%%%%%%%%%%%%%%%%%%%%%%%%%%%%%%%%%%%%%%%%%%%%%%%%%%
\begin{titlepage}
	\centering
	\vspace*{1cm}
	
	%— Smaller logo (20% of text width)
	\begin{figure}[H]
		\centering
		
\includegraphics[width=0.20\textwidth]{images/logo-fa.png}
	\end{figure}
	
	\vspace{1cm}
	
	{\fontsize{18}{28}\selectfont \textbf{\CourseTitle} }
	
	\vspace{3cm}
	
	{\fontsize{22}{28}\selectfont \textbf{\DocTitle} }
	
	\vspace{3cm}
	
	{\fontsize{14}{22}\selectfont \rl{نگارش:}}\\
	\vspace{0.5cm}
	{\fontsize{16}{22}\selectfont \textbf{ \Writers } }

	
	\vspace{1cm}
	
	{\fontsize{14}{22}\selectfont \rl{استاد درس:}}\\
	\vspace{0.5cm}
	{\fontsize{16}{22}\selectfont \textbf{ \FirstProfessor } }
	
	
	\vfill
	
	{\fontsize{14}{18}\selectfont \textbf{\Faculty} }\\
	\vspace{0.5cm}
	{\fontsize{14}{18}\selectfont \textbf{\University} }
	\vspace{0.5cm}
	
	\par
	
	{\fontsize{14}{18}\selectfont \textbf{\Semester} }
	
\end{titlepage}

%— Switch to our fancy header/footer from this page onward
\clearpage

%%%%%%%%%%%%%%%%%%%%%%%%%%%%%%%%%%%%%%%%%%%%%%%%%%%%%%%%%%%%
% === TABLE OF CONTENTS ===
%%%%%%%%%%%%%%%%%%%%%%%%%%%%%%%%%%%%%%%%%%%%%%%%%%%%%%%%%%%%
\tableofcontents
\clearpage
\pagestyle{aftertitle}

%%%%%%%%%%%%%%%%%%%%%%%%%%%%%%%%%%%%%%%%%%%%%%%%%%%%%%%%%%%%
% === MAIN CONTENT ===
%%%%%%%%%%%%%%%%%%%%%%%%%%%%%%%%%%%%%%%%%%%%%%%%%%%%%%%%%%%%
\chapter{\DocTitle}
\section{مقدمه}

برای کاربردی‌تر کردن برنامه‌های جاوایتان، نیاز است با \lr{GUI} آشنا بشوید.
\lr{GUI} که مخفف \lr{Graphical User Interface} (رابط گرافیکی کاربر) است،
در واقع همان پنجره‌ها و عناصر گرافیکی‌ است که توی سیستم عامل‌های مختلف برای
راحت‌تر شدن کار با برنامه‌ها استفاده می‌شود. چیزهایی مثل دکمه‌ها، برچسب‌ها،
جدول‌ها، فرم‌ها، فیلدهای ورودی، تصاویر و کلی عناصر دیگر که باعث راحت‌تر شدن
تعامل با برنامه می‌شوند.

با استفاده از \lr{GUI} در برنامه‌تان، کاربر می‌تواند بدون نیاز به نوشتن کد
یا وارد کردن دستی اطلاعات، فقط با یک کلیک، عملیات مورد نظرش را انجام
بدهد.

جاوا برای ساخت برنامه‌های گرافیکی، اول پکیج \lr{\texttt{AWT}} و بعد پکیج
\lr{\texttt{Swing}} را معرفی کرد. این دو، کلی شباهت و مقداری تفاوت با هم
دارند. در ادامه، اول نگاه کوتاهی به مفاهیم اشیا و کلاس‌ها می‌اندازیم و سپس
هم بخش‌های مهم \lr{\texttt{Swing}} و \lr{\texttt{AWT}} را با هم یاد
می‌گیریم.

\section{اولین برنامه گرافیکی خود را بنویسید}

کد زیر را در تابع \lr{main} قرار دهید. فعلا نگران جزئیات کد نباشید و
خودتان را گیج نکنید. به زودی قسمت‌به‌قسمت آن را توضیح خواهیم داد. :)

\begin{codebox}{java}
public class Main{    
   public static void main(String[] args) {  
      JFrame frame = new JFrame("Hello, java");  
      frame.setSize(300, 150);  
      JLabel label = new JLabel("Hello, Java!", JLabel.CENTER);  
      frame.add(label);  
      frame.setVisible(true);
   }
\end{codebox}

کد بالا را ران کنید. با همچین خروجی‌ای مواجه خواهید شد:

\begin{figure}[H]
\centering
\includegraphics[width=0.8\textwidth]{images/graphics-hello-world.png}
\end{figure}

\subsection{کلاس‌ها و آبجکت‌های برنامه \lr{Hello Java}}

حالا می‌خواهیم کدی که زده‌ایم را بررسی کنیم:

\begin{codebox}{java}
JFrame frame = new JFrame("Hello, java");
\end{codebox}

در اینجا از \lr{\texttt{JFrame}} استفاده کردیم که کلاسی از پکیج سویینگ
است که یک پنجره گرافیکی را نشان می‌دهد که بقیه‌ی عناصر گرافیکی روی آن سوار
می‌شوند. اینجا یک کلاس از این آبجکت ساختیم و اسم آن را \lr{frame}
گذاشتیم. ورودی (استرینگِ \lr{\texttt{"Hello, Java"}} ) مشخص می‌کند که چه
عنوانی برای پنجره‌ی ما قرار بگیرد. خودتان امتحان کنید و ببینید که وقتی به
این بخش ورودی ندهید، پنجره‌تان نیز تایتل ندارد.

\begin{codebox}{java}
frame.setSize(300, 150);
\end{codebox}

متد \lr{\texttt{setSize()}} توی کلاس \lr{\texttt{JFrame}} تعریف شده است
و برای هر آبجکتی که آن را فراخوانی می‌کند، بسته به ورودی‌ای که به آن داده
می‌شود، ابعاد فریم (همان پنجره) را مشخص می‌کند. (ورودی‌ها به ترتیب از چپ،
طول و عرض را بر حسب پیکسل مشخص می‌کنند).

\begin{codebox}{java}
JLabel label = new JLabel("Hello, Java!", JLabel.CENTER);
\end{codebox}

کلاس \lr{\texttt{JLabel}} مثل یک برچسب فیزیکی می‌ماند. آبجکت ساخته‌شده از
این کلاس، متن مورد نظرمان را در بخش مشخصی از فریم (در این مثال، به خاطر
استفاده از \lr{\texttt{CENTER}} ، در مرکز) قرار می‌دهد. در اینجا از یک
مفهوم کاملا شی‌گرا استفاده کردیم، چون از یک آبجکت برای قرار دادن متنمان
استفاده کردیم به جای اینکه خیلی ساده، یک متد را روی آبجکت \lr{frame}
فراخوانی کنیم تا متن را بنویسد. به این ترتیب یک ماهیت جداگانه به متنمان
دادیم تا بتواند ویژگی‌های مستقل خود را داشته باشد. بعد با فراخوانی متد
\lr{\texttt{add()}} روی آبجکت \lr{\texttt{frame}}، لیبل خود را به فریم
اضافه کردیم:

\begin{codebox}{java}
frame.add(label);
\end{codebox}

در مرحله آخر هم، فریممان و اجزای داخلش که کامل شده‌اند را نمایش دادیم:

\begin{codebox}{java}
frame.setVisible(true);
\end{codebox}

\section{آشنایی با \lr{AWT} و \lr{Swing}}

کلاس‌های \lr{\texttt{AWT}} امکانات زیادی را در اختیار برنامه‌نویس‌ها قرار
می‌دهند، اما یک نقطه‌ضعف بزرگ دارند؛ و آن این است که ظاهر برنامه‌هایی که با
\lr{\texttt{AWT}} طراحی می‌شوند، روی پلتفرم‌ها یا سیستم‌عامل‌های مختلف
می‌تواند تا حدی متفاوت باشد. علاوه بر این، نحوه‌ی تعامل کاربر با
برنامه‌هایی که با \lr{\texttt{AWT}} ساخته می‌شوند، در سیستم‌عامل‌های مختلف
ممکن است تفاوت کند. در واقع، وقتی یک برنامه را با \lr{\texttt{AWT}}
طراحی می‌کنید، این برنامه در هر سیستم‌عاملی که اجرا شود، از همان ابزارهای
گرافیکی پیش‌فرض آن سیستم‌عامل برای نمایش رابط کاربری استفاده می‌کند. این
موضوع باعث می‌شود که برنامه شما در هر سیستم‌عامل، یک ظاهر متفاوت داشته
باشد. برای حل این مشکل، شرکت اوراکل \lr{\texttt{Swing}} را به زبان جاوا
اضافه کرد.

کتابخانه \lr{\texttt{Swing}} در جاوا یک مجموعه از کلاس‌های از پیش تعریف
شده است که امکانات زیادی برای ساخت \lr{GUI} در اختیار شما می‌گذارد.
کلاس‌های \lr{\texttt{Swing}} نسبت به \lr{\texttt{AWT}} انعطاف‌پذیری بیشتری
دارند و ویژگی‌هایشان نیز پیشرفته‌تر است. مثلا جدول، فرم، لیست، اسکرول و
بسیاری المان دیگر را در \lr{\texttt{Swing}} می‌توانید راحت‌تر و با کنترل
بیشتری طراحی کنید.

قبل از اینکه وارد جزئیات \lr{\texttt{Swing}} شویم، یک یادآوری کوتاه از
شی‌گرایی لازم داریم؛ چون تقریبا تمام چیزی که در \lr{GUI} جاوا می‌سازیم،
دقیقا با همین مفاهیم جلو می‌رود. در برنامه‌‌های گرافیکی، هر چیزی که روی
صفحه می‌بینید (مثل دکمه، برچسب، پنل و حتی خود پنجره) یک آبجکت از یک کلاس
مشخص است؛ یعنی به جای اینکه «مستقیم روی صفحه نقاشی کنیم»، داریم مجموعه‌ای
از اشیاء می‌سازیم، به آن‌ها ویژگی می‌دهیم (مثل اندازه، رنگ، متن و موقعیت) و
بعد آن‌ها را کنار هم قرار می‌دهیم. پس همان مفاهیمی که در شی‌گرایی گفتیم مثل
ساختن آبجکت با \lr{\texttt{new}}، صدا زدن متدها روی آبجکت، و جدا کردن
مسئولیت‌ها بین کلاس‌های مختلف، اینجا کاملا عملی و ملموس می‌شوند. از طرفی،
خود کلاس‌های \lr{GUI} جاوا هم با همان منطق شی‌گرایی طراحی شده‌اند: یعنی
داخل یک سلسله‌مراتب قرار دارند و از هم ارث‌بری می‌کنند.

ارتباط \lr{\texttt{AWT}} و \lr{\texttt{Swing}} و ساختار سلسله‌مراتبی
آن‌ها:

\begin{figure}[H]
\centering
\includegraphics[width=0.8\textwidth]{images/graphics-swing-hierarchy.png}
\end{figure}

طبق نمودار، سلسله‌مراتب از\lr{\texttt{Object}} شروع می‌شود و به
\lr{\texttt{Component}} می‌رسد. \lr{\texttt{Component}} یک کلاس در
\lr{\texttt{AWT}} است و هر چیزی که قرار است روی صفحه نمایش داده شود، در
نهایت از این کلاس ارث می‌برد. بعد از \lr{\texttt{Component}}، دو شاخه‌ی
مهم می‌بینیم: شاخه‌ی \lr{\texttt{Container}} و شاخه‌ی
\lr{\texttt{JComponent}}. کلاس \lr{\texttt{Container}} (مثل
\lr{\texttt{Window}} و \lr{\texttt{Panel}}) خودش یک
\lr{\texttt{Component}} است، با این تفاوت که نقش ظرف را دارد؛ یعنی
می‌تواند چند \lr{\texttt{Component}} دیگر را داخل خودش نگه دارد و به همین
خاطر برای ساختاردهی صفحه (مثل پنجره‌ها و پنل‌های پایه در \lr{AWT)} خیلی
کاربردی است. در سمت دیگر، \lr{\texttt{JComponent}} هم زیرکلاس
\lr{\texttt{Component}} است، اما بیشتر به دنیای \lr{\texttt{Swing}}
مربوط می‌شود و پایه‌ی اکثر ویجت‌های \lr{Swing} به حساب می‌آید؛ یعنی کلاس‌هایی
مثل \lr{\texttt{JLabel}}، \lr{\texttt{JList}}، \lr{\texttt{JTable}} و
\lr{\texttt{JButton}} از آن مشتق می‌شوند. پس می‌توان گفت که هر دو شاخه در
نهایت \lr{component} بودن را به ارث می‌برند، اما \lr{\texttt{Container}}
بیشتر برای نگه‌داری و سازماندهی اجزا است و \lr{\texttt{JComponent}} بیشتر
برای ساختن ویجت‌های آماده‌ی \lr{Swing .} در ادامه، بخش‌‌های
\lr{\texttt{Window}} ،\lr{\texttt{Panel}} و \lr{\texttt{JComponent}} را
به همراه زیرکلاس‌های هر بخش بررسی می‌کنیم.

\subsection{\lr{Window}}

پنجره‌ای است که هنگام اجرای برنامه باز می‌شود و می‌تواند به دو حالت باشد:
\lr{\texttt{dialog}} (حالت پاپ‌آپی که در برنامه‌ها می‌بینید) و
\lr{\texttt{Frame}} (قابی که پنجره هر برنامه با اجرا شدن خودش، باز
می‌کند). بنابراین برای ایجاد یک رابط گرافیکی، به یک \lr{\texttt{Frame}}
نیاز داریم و \lr{\texttt{JFrame}} در جاوا این کار را برای ما انجام
می‌دهد. \lr{\texttt{JFrame}} یک \lr{Container} یا ظرف است که این قابلیت
را دارد که Component‌ هایی مثل \lr{\texttt{JButton}} ،
\lr{\texttt{JTextArea}} ، \lr{\texttt{JPanel}} و \ldots{} را در خودش جای
دهد.

در ابتدا باید کتابخانه \lr{\texttt{swing}} را برای ایجاد رابط گرافیکی به
عنوان پیشنیاز کد ایمپورت کنیم. سپس یک آبجکت از کلاس ‍‍\lr{\texttt{JFrame}}
می‌‌سازیم (اینجا مانند ابتدای اکیومنت نام آن را \lr{\texttt{frame}}
گذاشتیم). بعد از آن که آبجکت را ساختیم، برای نمایش آن باید مقدار
\lr{\texttt{visibility}} را \lr{\texttt{true}} قرار بدهیم. نمونه کد و
نتیجه‌ی اجرای آن را در ادامه می‌بینید. سام و نسترن قصد دارند برای تیم
تدریسیاری، رابط کاربریِ فرمی را طراحی کنند که دانشجویانی که این ترم درس
\lr{AP} دارند، بتوانند اطلاعاتشان را توسط آن ثبت کنند. بنابراین، نام
پنجره ایجاد شده را \lr{Student Registration Form} می‌گذاریم.

\begin{codebox}{java}
import javax.swing.*;  
public class Main{    
   public static void main(String[] args) {  
      JFrame frame = new JFrame("Student Registration Form");  
      frame.setVisible(true);  
   }
\end{codebox}

خروجی کد :

\begin{figure}[H]
\centering
\includegraphics[width=0.8\textwidth]{images/graphics-swing-1.png}
\end{figure}

همانطور که می‌بینید، خروجی فقط یک پنجره‌ی خالی است. هیچ چیزی داخلش نیست و
عنوانش هم دیده نمی‌شود. وقتی اندازه‌ی پنجره را دستی تغییر بدهید، می‌بینید
که قاب باز شده به این شکل نشان داده می‌شود:

\begin{figure}[H]
\centering
\includegraphics[width=0.8\textwidth]{images/graphics-swing-2.png}
\end{figure}

حالا می‌توانیم ببینیم که پنجره داده شده عنوان دارد ولی هنوز محتوایی
ندارد. در واقع، یک آبجکت از کلاس \lr{\texttt{JFrame}} مثل یک قاب یا
تابلوی نقاشی است که هنوز چیزی روی آن کشیده نشده. قبل از اینکه اجزای
دیگری به آن اضافه کنیم، بیایید کمی با همین قاب خالی کار کنیم.

مثلا می‌خواهیم اندازه‌ی پنجره‌ی باز شده بعد از اجرا را کمی بزرگتر کنیم. اول
باید ابعاد مورد نظر را بسازیم، اما برای این کار نیاز به اضافه کردن یک
کتابخانه دیگر هم داریم که باید آن را هم به ابتدای کد اضافه کنیم:

\begin{codebox}{java}
import java.awt.*;
\end{codebox}

سپس، ابعاد را به این صورت می‌سازیم:

\begin{codebox}{java}
Dimension frameSize = new Dimension(1024,720);

// Create a Dimension object to set the window size (width:1024, height:720)
\end{codebox}

سپس این ابعاد را به \lr{\texttt{frame}} می‌دهیم تا روی آن اعمال شود. این
کار را به صورت زیر انجام می‌دهیم:

\begin{codebox}{java}
frame.setSize(frameSize);
\end{codebox}

کد آن به صورت زیر است:

\begin{codebox}{java}
import javax.swing.*;
import java.awt.*;
public class Main{
   public static void main(String[] args) {  
      JFrame frame = new JFrame("Student Registration Form");  
      Dimension frameSize = new Dimension(1024,720);  
      frame.setSize(frameSize);  
      frame.setVisible(true);
   }
}
\end{codebox}

خروجی کد:

\begin{figure}[H]
\centering
\includegraphics[width=0.8\textwidth]{images/graphics-swing-3.png}
\end{figure}

دقت کنید که وقتی پنجره را می‌بندید، کد شما همچنان ادامه پیدا می‌کند و
متوقف نمی‌شود. برای این که این رفتار را تغییر دهید، می‌توانید حالت عملیات
بستن پنجره را با کد زیر تغییر دهید. (گزینه‌ای که در کد زیر تحت عنوان
\lr{\texttt{EXIT\_ON\_CLOSE}} استفاده شده باعث می‌شود کد به طور کامل با
بسته شدن پنجره متوقف شود. حالت‌های دیگری هم علاوه بر
\lr{\texttt{EXIT\_ON\_CLOSE}} هست که می‌توانید از آن‌ها استفاده کنید.)

\begin{codebox}{java}
frame.setDefaultCloseOperation(JFrame.EXIT_ON_CLOSE);
\end{codebox}

کد نهایی ما به صورت زیر است:

\begin{codebox}{java}
import javax.swing.*;
import java.awt.*;
public class Main{
   public static void main(String[] args) {  
      JFrame frame = new JFrame("Student Registration Form");  
      Dimension frameSize = new Dimension(1024,720);  
      frame.setSize(frameSize);  
      frame.setVisible(true);  
      frame.setDefaultCloseOperation(JFrame.EXIT_ON_CLOSE);  
   }  
}
\end{codebox}

توجه کنید که برای مشخص کردن ابعاد پنجره‌تان، می‌توانید از کد زیر هم
استفاده کنید:

\begin{codebox}{java}
frame.setSize(500, 500);
\end{codebox}

در ادامه، به اجزای مختلفی مي‌پردازیم که مي‌توانید به این قاب اضافه کنید.

\subsection{\lr{Panel}}

\lr{JPanel} ساده‌ترین کلاس در بین اجزای گرافیکی در جاواست.
\lr{\texttt{JPanel}} فضایی را در برنامه ایجاد می‌کند که می‌توانید هر جزء
گرافیکی را به آن اضافه کنید. می‌توانید \lr{\texttt{JPanel}} را مثل بوم
نقاشی تصور کنید که اجزای گرافیکی شما مثل اجزای نقاشی روی آن قرار
می‌گیرند.

حالا ممکن است این سؤال پیش بیاید که آیا \lr{\texttt{JFrame}} و
\lr{\texttt{JPanel}} مثل هم نیستند؟ چون در \lr{\texttt{JFrame}} هم اجزای
گرافیکی را به آن اضافه می‌کردیم. باید بگویم که \lr{\texttt{JFrame}} فقط
یک پنجره معمولی است که در برنامه‌های کاربردی از آن استفاده می‌کنیم. اما
\lr{\texttt{JPanel}} با امکاناتی که دارد برای سازماندهی اجزای گرافیکی در
جاوا خیلی بهتر و مناسب‌تر است. در واقع، \lr{\texttt{JPanel}} جزئی است که
روی این قاب ( \lr{\texttt{JFrame}} ) قرار می‌گیرد و تقریبا خودش می‌تواند
مثل یک \lr{mini frame} برای اجزای دیگر مثل دکمه‌ها، لیبل‌ها و\ldots{}
باشد.

معمولا در یک برنامه، از یک فریم اصلی استفاده می‌شود و پنل‌ها روی این فریم
قرار می‌گیرند. طی تعامل کاربر با برنامه (مثلا از طریق دکمه‌ها)، این پنل‌ها
می‌توانند جابجا شوند. مثلا وقتی یک دکمه را می‌زنید، از یک صفحه به صفحه
دیگر می‌روید. جلوتر برای این مورد مثال خواهیم زد.

در ادامه می‌خواهیم پنلی به رنگ آبی روی قسمتی از پنجره‌ی خود ایجاد کنیم.
اما برای آن که بتوانید پنل‌ها را آزادانه روی قاب قرار بدهید، ابتدا باید
برای قاب خود چیدمانی خالی تعیین کنید، در غیر این صورت ممکن است چیدمان‌های
پیش‌فرض جاوا جلوی شما را بگیرند (یک چیدمان‌ یا \lr{\texttt{Layout}}
مجموعه‌ای از قواعد یا به اصطلاح استایل قرار گرفتن اجزای گرافیکی در یک
پنجره است. در ادامه به طور مختصر به انواع چیدمان‌‌ها می‌پردازیم). برای این
کار، ورودی متد زیر را \lr{\texttt{null}} قرار می‌دهیم. توجه کنید که در کد
زیر، \lr{\texttt{null}} یعنی هیچ‌گونه \lr{\texttt{Layout Manager}} برای
فریم تنظیم نشده. یعنی شما باید خودتان موقعیت و ابعاد اجزای رابط کاربری
را به صورت دستی مشخص کنید. در واقع، در اینجا وقتی از \lr{\texttt{null}}
استفاده می‌کنید، ترتیب قرارگیری اجزا و اندازه‌های آن‌ها به شما واگذار می‌شود
و دیگر \lr{\texttt{Java}} به طور خودکار آن‌ها را بر اساس یک طرح‌بندی خاص
تنظیم نمی‌کند.

\begin{codebox}{java}
frame.setLayout(null);
\end{codebox}

سپس یک آبجکت جدید از روی این کلاس \lr{JPanel} ایجاد می‌کنیم:

\begin{codebox}{java}
JPanel mainPanel = new JPanel();
\end{codebox}

به کمک کد زیر می‌توانیم اندازه و موقعیت پنلمان را مشخص کنیم. (اندازه‌ی پنل
می‌تواند با اندازه‌ی پنجره هم یکسان باشد.):

\begin{codebox}{java}
mainPanel.setBounds(40, 35, 400, 400);
\end{codebox}

نکته‌ای که اهمیت دارد این است که در تنظیم مرزهای پنل، مقداردهی‌های داده
شده به صورت زیر تحلیل می‌شوند:

دو مولفه اول (35, 40) به ترتیب طول و عرض نقطه‌ای \lr{(x, y)} هستند که پنل
ما از آن نقطه شروع می‌شود. مولفه‌های سوم و چهارم نیز به ترتیب عرض و ارتفاع
پنل هستند. به مثال زیر در این راستا توجه کنید:

\begin{codebox}{java}
panel.setBounds(0, 0, 512, 720);
\end{codebox}

\begin{itemize}
\tightlist
\item
  \textbf{مولفه‌ی اول و دوم}: مقدار \lr{\texttt{0}} سمت چپ مکان شروع پنل
  روی بردار \lr{\texttt{x}}‌ها است و مقدار \lr{\texttt{0}} سمت راست مکان
  شروع پنل روی محور \lr{\texttt{y}}‌ها است (به شکل زیر دقت کنید. محورها
  روی قاب به این شکل هستند).
\end{itemize}

\begin{figure}[H]
\centering
\includegraphics[width=0.8\textwidth]{images/graphics-swing-4.png}
\end{figure}

\begin{itemize}
\tightlist
\item
  مولفه سوم و چهارم :مقادیر \lr{\texttt{512}} و \lr{\texttt{720}} به
  ترتیب عرض و ارتفاع پنل هستند.
\item
  از متد \lr{\texttt{setBounds()}} زمانی که \lr{\texttt{layout}}ای
  نداریم (آن را \lr{\texttt{null}} گذاشته‌ایم) استفاده می‌کنیم.
\end{itemize}

برای تنظیم رنگ پنل، کلاس \lr{Color} چند رنگ پیش‌فرض دارد (با نوشتن
\lr{\texttt{Color.}} می‌توانید آن‌ها را ببینید)، اما راستش را بخواهید،
معمولا رنگ‌های جذابی نیستند. راه بهتر این است که یک شی جدید از \lr{Color}
بسازید و کد هگزادسیمال رنگ دلخواه را به آن بدهید؛ مثلا من از کد رنگ
\lr{baby blue} استفاده کرده‌ام.

کد این رنگ \lr{\texttt{\#89CFF0}} است، اما اگر از مبانی برنامه‌نویسی
یادتان باشد، به جای \lr{\texttt{\#}} ، با \lr{\texttt{0x}} در ابتدای عدد
مشخص می‌کنیم که عدد ما هگزادسیمال است (\lr{\texttt{0x89CFF0}}). در کد زیر
رنگ پنل را مشخص کردیم. از این فرایند برای تعیین رنگ هر کامپوننت دیگری
نیز می‌توانید استفاده کنید.

\begin{codebox}{java}
mainPanel.setBackground(new Color(0x89CFF0));
\end{codebox}

در انتها باید به صورت زیر پنلمان را به \lr{\texttt{frame}} اضافه کنیم:

\begin{codebox}{java}
frame.add(mainPanel);
\end{codebox}

کد نهایی ما به صورت زیر خواهد بود:

\begin{codebox}{java}
import javax.swing.*;
import java.awt.*;

public class Main{
   public static void main(String[] args) {  
  
      JFrame frame = new JFrame("Student Registration Form");  
      frame.setSize(500, 500);  
      frame.setDefaultCloseOperation(JFrame.EXIT_ON_CLOSE);  
      frame.setLayout(null);  
  
      JPanel mainPanel = new JPanel();  
      mainPanel.setBounds(40, 35, 400, 400);   
      mainPanel.setBackground(new Color(0x89CFF0));

      frame.add(mainPanel);  
      frame.setVisible(true);  
   }  
}
\end{codebox}

خروجی کد فوق:

\begin{figure}[H]
\centering
\includegraphics[width=0.8\textwidth]{images/graphics-swing-5.png}
\end{figure}

\subsection{\lr{Labels}}

برچسب \lr{(Label)} یک آبجکت از کلاس \lr{\texttt{JLabel}} است که برای
نمایش متن یا عکس در فضای رابط کاربری استفاده می‌شود و فقط برای خواندن
قابل دسترس است (یعنی کاربر نمی‌تواند مستقیما آن را تغییر دهد). به طور
کلی، از کلاس \lr{\texttt{JLabel}} برای نمایش متن یا تصویر در برنامه‌های
گرافیکی استفاده می‌کنیم.

باز هم مثل قبل، برای استفاده از کلاس \lr{JLabel} باید ابتدا یک آبجکت
جدید از این کلاس بسازیم و یک نام برای آن انتخاب کنیم. از این به بعد
می‌توانیم با استفاده از نام این آبجکت، متدهای کلاس \lr{\texttt{JLabel}}
را فراخوانی کنیم.

در کد زیر، متن دلخواه را با استفاده از \lr{\texttt{JLabel}} نمایش می‌دهیم
و با استفاده از \lr{\texttt{setBounds}} موقعیت و اندازه برچسب‌ها را مشخص
می‌کنیم:

\begin{codebox}{java}
JLabel nameLabel = new JLabel("Student name:");  
nameLabel.setBounds(10, 10, 300, 50);  
JLabel numberLabel = new JLabel("Student number:");  
numberLabel.setBounds(10, 90, 300, 50);  
JLabel emailLabel = new JLabel("Email:");  
emailLabel.setBounds(10, 170, 300, 50);
\end{codebox}

در نهایت با کد زیر برچسب‌هایمان را به پنل اضافه می کنیم:

\begin{codebox}{java}
mainPanel.add(nameLabel);  
mainPanel.add(numberLabel);  
mainPanel.add(emailLabel);
\end{codebox}

کد نهایی به صورت زیر است:

\begin{codebox}{java}
import javax.swing.*;
import java.awt.*;

public class Main{
   public static void main(String[] args) {  
      JFrame frame = new JFrame("Student Registration Form");  
      frame.setSize(500, 500);  
      frame.setDefaultCloseOperation(JFrame.EXIT_ON_CLOSE);
      frame.setLayout(null);  
          
      JPanel mainPanel = new JPanel();  
      mainPanel.setBounds(40, 35, 400, 400);  
      mainPanel.setBackground(new Color(0x89CFF0));  
      mainPanel.setLayout(null);   
          
      JLabel nameLabel = new JLabel("Student name:");  
      nameLabel.setBounds(10, 10, 300, 50);  
      JLabel numberLabel = new JLabel("Student number:");  
      numberLabel.setBounds(10, 90, 300, 50);  
      JLabel emailLabel = new JLabel("Email:");  
      emailLabel.setBounds(10, 170, 300, 50);  
          
      mainPanel.add(nameLabel);  
      mainPanel.add(numberLabel);  
      mainPanel.add(emailLabel);  
          
      frame.add(mainPanel);  
      frame.setVisible(true);  
   }  
}
\end{codebox}

خروجی کد:

\begin{figure}[H]
\centering
\includegraphics[width=0.8\textwidth]{images/graphics-swing-6.png}
\end{figure}

\subsection{\lr{TextField}}

کلاس ‍\lr{\texttt{JTextField}} شامل اجزای گرافیکی «متنی» است که وقتی از
آن آبجکتی بسازید، به شما این امکان را می‌دهد که یک خط متن را ویرایش کنید.
برای اینکه بهتر متوجه شوید، می‌توانید برنامه \lr{Notepad} را باز کنید و
از منوی \lr{Edit} گزینه \lr{Find} را انتخاب کنید.

\begin{figure}[H]
\centering
\includegraphics[width=0.8\textwidth]{images/graphics-swing-7.png}
\end{figure}

در تصویر بالا، ما قصد داریم کلمه‌ی «جاوا» را در میان متن موجود در برنامه
\lr{Notepad} پیدا کنیم. فیلدی که در بخش \lr{Find} برای وارد کردن کلمه‌ی
مورد نظر استفاده می‌کنیم، همون جزء گرافیکی \lr{\texttt{JTextField}} است
که با رنگ قرمز مشخص کردیم. احتمالا با این نوع اجزای گرافیکی زیاد برخورد
کرده‌اید، مثلا وقتی که در فرم‌های ثبت‌نام، اطلاعات کاربری را وارد می‌کنید.
متنی که در این اجزای گرافیکی وارد می‌کنیم ویژگی‌های زیر را دارد:

\begin{itemize}
\tightlist
\item
  به صورت یک خط یا سطر است.
\item
  بر خلاف \lr{\texttt{JLabel}} که فقط برای نمایش متن استفاده می‌شود، این
  فیلدها توسط کاربر نیز قابل ویرایش هستند، یعنی شما می‌توانید متن قبلی را
  با متن جدید عوض کنید.
\end{itemize}

حالا می‌خواهیم سه \lr{\texttt{TextField}} برای سه لیبلی که قبلا ساختیم،
ایجاد کنیم:

\begin{codebox}{java}
JTextField nameTextField = new JTextField();  
nameTextField.setBounds(170, 10, 200, 40);  
JTextField numberTextField = new JTextField();  
numberTextField.setBounds(170, 90, 200, 40);  
JTextField emailTextField = new JTextField();  
emailTextField.setBounds(170, 170, 200, 40);
\end{codebox}

در انتها با کد زیر، \lr{\texttt{TextField}}‌هایمان را به پنل اضافه
می‌کنیم:

\begin{codebox}{java}
mainPanel.add(nameTextField);  
mainPanel.add(numberTextField);  
mainPanel.add(emailTextField);
\end{codebox}

خروجی کد:

\begin{figure}[H]
\centering
\includegraphics[width=0.5\textwidth]{images/graphics-swing-8.png}
\end{figure}

دقت شود که ما متن پیش‌فرضی برای \lr{\texttt{textfield}}‌هایمان نگذاشته
بودیم و متن موجود در فیلدهای مثال بالا، مثلا توسط کاربر وارد شده‌اند.

\subsection{نمایش تصویر با استفاده از کلاس \texttt{JLabel}}

همانطور که قبلا گفتیم، در جاوا می‌توانیم برای نمایش تصاویر از کلاس
\lr{\texttt{JLabel}} استفاده کنیم. این کلاس به ما این امکان را می‌دهد که
تصاویر را در صفحه نمایش برنامه نشان دهیم.

برای شروع، باید یک شی از کلاس \lr{\texttt{JLabel}} بسازیم و تصویر مورد
نظر را به آن اختصاص دهیم. برای این کار می‌توانیم از کلاس
\lr{\texttt{ImageIcon}} استفاده کنیم تا تصویر را بخوانیم و بعد آن را به
آبجکت ساخته شده از کلاس \lr{\texttt{JLabel}} اختصاص دهیم.
\lr{\texttt{ImageIcon}}، یک کلاس در سویینگ است که به ما کمک می‌کند تصاویر
را روی کامپوننت‌هایی مثل \lr{\texttt{JLabel}} و \lr{\texttt{JButton}}
بارگذاری کنیم و نمایش دهیم. مثلا برای نمایش تصویر با نام
\lr{\texttt{"imagee.png"}} که در پوشه \lr{\texttt{src}} برنامه قرار
دارد، به صورت زیر عمل می‌کنیم:

\begin{codebox}{java}
ImageIcon imageIcon = new ImageIcon("src/imagee.png");
\end{codebox}

\begin{figure}[H]
\centering
\includegraphics[width=0.8\textwidth]{images/graphics-swing-9.png}
\end{figure}

\begin{figure}[H]
\centering
\includegraphics[width=0.8\textwidth]{images/graphics-swing-10.png}
\end{figure}

مسیر دقیق فایل را به صورت بالا در کدمان مشخص می‌کنیم و با استفاده از کلاس
\lr{\texttt{ImageIcon}} تصویر \lr{\texttt{"imagee.png"}} را می‌خوانیم. در
ادامه برای تغییر اندازه تصویر \lr{(scaling)} از کد زیر استفاده می‌کنیم:

\begin{codebox}{java}
Image scaledImage = imageIcon.getImage().getScaledInstance(150, 150, Image.SCALE_SMOOTH);
\end{codebox}

\begin{itemize}
\tightlist
\item
  متد \lr{\texttt{getImage()}} از کلاس \lr{\texttt{ImageIcon}} تصویر خام
  \lr{(Image)} را که داخل آبجکت \lr{\texttt{ImageIcon}} است، برمی‌گرداند.
  می‌توانید از این متد برای کارهای گرافیکی یا تغییر اندازه‌ی تصویر استفاده
  کنید.
\item
  متد \lr{\texttt{getScaledInstance()}} از کلاس \lr{\texttt{Image}} برای
  تغییر اندازه‌ی تصویر به کار می‌رود. این متد سه تا ورودی می‌گیرد:
\item
  عرض تصویر (\lr{\texttt{width}}): مقدار عددی عرض جدید تصویر که در اینجا
  \lr{\texttt{150}} پیکسل تعیین می کنیم.
\item
  ارتفاع تصویر (\lr{\texttt{height}}): مقدار عددی ارتفاع جدید تصویر که
  اینجا هم \lr{\texttt{150}} پیکسل تعیین می‌کنیم.
\item
  راهنماهای مقیاس‌دهی (\lr{\texttt{hints}}): روشی که می‌خواهیم برای تغییر
  اندازه‌ی تصویر استفاده کنیم. این آرگومان الگوریتم مقیاس‌دهی را مشخص
  می‌کند. در اینجا از \lr{\texttt{Image.SCALE\_SMOOTH}} استفاده کردیم.
\end{itemize}

مقدار \lr{\texttt{Image.SCALE\_SMOOTH}} یک ثابت از کلاس
\lr{\texttt{Image}} است که الگوریتمی برای مقیاس‌دهی تصویر ارائه می‌دهد.
این روش باعث می‌شود کیفیت تصویر موقع تغییر اندازه حفظ شود. سرعتش شاید کمی
پایین‌تر باشد، اما نتیجه‌ی نهایی کیفیت بالاتری دارد. بعضی مقادیر دیگر برای
\lr{\texttt{hints}} عبارتند از:

\begin{itemize}
\tightlist
\item
  \lr{\texttt{Image.SCALE\_FAST}} : سرعت بالاتر با کیفیت پایین‌تر.
\item
  \lr{\texttt{Image.SCALE\_DEFAULT}} : استفاده از تنظیمات پیش‌فرض سیستم.
\item
  \lr{\texttt{Image.SCALE\_REPLICATE}} : مقیاس‌دهی ساده که ممکن است کیفیت
  خوبی نداشته باشد.
\end{itemize}

مثالی برای اینکه این موضوع را بهتر متوجه شوید: فرض کنید تصویری دارید که
اندازه‌اش \lr{\texttt{750*500}} پیکسل است و می‌خواهید آن را کوچک کنید تا
در فضای \lr{\texttt{150*100}} پیکسلی جا شود. این خط دقیقا این کار را
انجام می‌دهد، بدون اینکه کیفیت تصویر خیلی افت کند.

حالا تصویر را به شی \lr{\texttt{JLabel}} اختصاص می‌دهیم:

\begin{codebox}{java}
JLabel imageLabel = new JLabel(new ImageIcon(scaledImage));
\end{codebox}

سپس به کمک \lr{\texttt{setBounds}} موقعیت تصویر را در صفحه مشخص می کنیم:

\begin{codebox}{java}
imageLabel.setBounds(30, 230, 100, 150);
\end{codebox}

و در انتها، تصویر را به پنلمان اضافه می کنیم:

\begin{codebox}{java}
mainPanel.add(imageLabel);
\end{codebox}

خروجی کد:

\begin{figure}[H]
\centering
\includegraphics[width=0.8\textwidth]{images/graphics-swing-11.png}
\end{figure}

\subsection{\lr{Button}}

همانطور که قبلا گفتیم، یکی از روش‌هایی که کاربر می‌تواند با برنامه تعامل
داشته باشد، کلیک کردن روی دکمه‌هاست. پس هر کدام از دکمه‌ها می‌توانند یک کار
خاصی را انجام دهند. اشیاء ساخته شده از کلاس \lr{\texttt{JButton}}
دکمه‌هایی هستند که وقتی کاربر روی آن‌ها کلیک می‌کند، یک کار مشخص را انجام
می‌دهند.

در کد زیر که دو دکمه تعریف کردیم، می‌توانید متن دلخواه را روی دکمه
بگذارید و با استفاده از \lr{\texttt{setBounds}} نیز موقعیت و اندازه‌ی
دکمه را مشخص کنید.

\begin{codebox}{java}
JButton submitButton = new JButton("Submit");  
submitButton.setBounds(10, 10, 300, 50);

JButton okButton = new JButton("OK");  
okButton.setBounds(150, 330, 100, 50);
\end{codebox}

بعد از اینکه دکمه را تعریف کردید، می‌توانید با استفاده از روش زیر آن را
فعال کنید. در این برنامه می‌خواهیم وقتی دکمه زده شد، پیامی مبتنی بر ثبت
شدن اطلاعات را نشان دهد. برای این کار دو حالت مختلف را با دو مثال بررسی
می‌کنیم.

اولا نیاز داریم که در ابتدای کدمان، پکیج زیر را ایمپورت کنیم:

\begin{codebox}{java}
Import java.awt.event.*;
\end{codebox}

حالا، کد زیر را به برنامه اضافه می‌کنیم:

\begin{codebox}{java}
submitButton.addActionListener(new ActionListener() {
   @Override
   public void actionPerformed(ActionEvent e) {  
      // Button functionality
   }
});
\end{codebox}

در این قطعه کد، برای دکمه، یک \lr{\texttt{ActionListener}} تعریف می‌کنیم
تا هر زمان کاربری روی دکمه کلیک کرد، متد \lr{\texttt{actionPerformed}}
اجرا شود. از جلسه‌ی شی‌ءگرایی احتمالا به یاد دارید که خیلی وقت‌ها ما با نوع
کلی کار می‌کنیم، نه با نوع دقیق. اینجا هم دقیقا به همین صورت است: متد
\lr{\texttt{addActionListener}} از ما یک ورودی از نوع
\lr{\texttt{ActionListener}} می‌خواهد اما برایش مهم نیست دقیقا چه کلاسی
باشد؛ فقط کافی است اینترفیس \lr{\texttt{ActionListener}} را پیاده‌سازی
کرده باشد (پیاده‌سازی متد \lr{\texttt{actionPerformed}} را داشته باشد).
به عبارت دیگر، ما یک «شنونده» می‌دهیم، ولی این شنونده می‌تواند هر
پیاده‌سازی‌ای داشته باشد و هر بار که کلیک اتفاق می‌افتد، همان نسخه‌ای که ما
نوشته‌ایم اجرا می‌شود. در قسمت کامنت‌ شده، کد مربوط به کارهایی را می‌نویسیم
که می‌خواهیم بعد از فشردن دکمه اجرا شود (متد
\lr{\texttt{actionPerformed}} را \lr{override} می‌کنیم). پس شکل فعال‌سازی
دکمه معمولا ثابت است، اما کاری که انجام می‌شود می‌تواند خیلی متنوع باشد و
بسته به نیازمان عوض شود که دقیقا یک نمونه کاربردی از \lr{Polymorphism}
است.

برای اضافه کردن شنونده به دکمه‌ها معمولا همین رویکرد بالا را استفاده
می‌کنیم: به‌جای اینکه یک کلاس جداگانه بسازیم که
\lr{\texttt{ActionListener}} را پیاده‌سازی کرده باشد، همان‌جا
\lr{(inline)} یک نمونه‌ی ناشناس می‌سازیم و متد
\lr{\texttt{actionPerformed}} را \lr{override} می‌کنیم تا رفتار دکمه را
مشخص کنیم. دلیلش هم روشن است: خیلی وقت‌ها دکمه‌های مختلف قرار نیست دقیقا
یک کار یکسان انجام دهند؛ اگر برای هر دکمه یک کلاس جدا می‌ساختیم، تعداد
کلاس‌های \lr{\texttt{ActionListener}} خیلی زیاد و مدیریت کد سخت می‌شد.
البته اگر واقعا تعداد زیادی دکمه با رفتار کاملا مشابه داشته باشیم، آن‌وقت
ساختن یک کلاس جداگانه و استفاده‌ی مجدد از آن می‌تواند انتخاب بهتری باشد.

همانطور که گفتیم، می‌خواهیم با فشردن دکمه، یک پیغام نشان دهیم که
نشان‌دهنده ثبت شدن اطلاعات باشد.

\textbf{مثال یک}

یک پنل جدید می‌سازیم و به آن یک لیبل جدید و دکمه اختصاص می‌دهیم. هدف این
است که با زدن دکمه، پنل‌ها روی فریم جایگزین شوند. برای پنل دوم، از یکی از
رنگ‌های پیش‌فرض کلاس \lr{Color} استفاده می‌کنیم:

\begin{codebox}{java}
JPanel successPanel = new JPanel();  
successPanel.setBounds(40, 35, 400, 400);   
successPanel.setBackground(Color.lightGray);  
successPanel.setLayout(null);

JLabel infoLabel = new JLabel("Information Submitted successfully");  
infoLabel.setBounds(100, 170, 350, 50);

successPanel.add(infoLabel);
\end{codebox}

دکمه‌ی \lr{\texttt{submitButton}} را به پنل اول و دکمه‌ی
\lr{\texttt{okButton}} را به پنل دوم اضافه می‌کنیم:

\begin{codebox}{java}
mainPanel.add(submitButton);  
successPanel.add(okButton);
\end{codebox}

حالا، \lr{\texttt{actionListener}} را برای هر دو دکمه پیاده سازی می‌کنیم،
به طوری که با زدن دکمه‌ها‌، بین پنل‌ها جابه‌جا شویم:

\begin{codebox}{java}
submitButton.addActionListener(new ActionListener() {
   @Override
   public void actionPerformed(ActionEvent e) {
      frame.getContentPane().removeAll();

      frame.add(successPanel);

      frame.revalidate();

      frame.repaint();

   }
});

okButton.addActionListener(new ActionListener() {
   @Override
   public void actionPerformed(ActionEvent e) {
      frame.getContentPane().removeAll();

      frame.add(mainPanel);

      frame.revalidate();

      frame.repaint();

   }
});
\end{codebox}

متد \lr{\texttt{getContentPane().removeAll()}} تمام اجزا و کامپوننت‌های
داخل فریم را حذف می‌کند. این کار باعث جلوگیری از باقی ماندن اجزای قبلی
هنگام تغییر پنل می‌شود. متد \lr{\texttt{revalidate()}} چیدمان
(\lr{\texttt{layout}}) را به‌روزرسانی می‌کند و متد \lr{\texttt{repaint()}}
کل فریم را دوباره رسم می‌کند تا تغییرات به درستی نشان داده شوند.

خروجی کد بعد از کلیک بر روی دکمه \lr{\texttt{Submit}}:

\begin{figure}[H]
\centering
\includegraphics[width=0.8\textwidth]{images/graphics-swing-12.png}
\end{figure}

با کلیک بر روی دکمه \lr{OK} ، به صفحه‌ی اصلی برمی‌گردیم.

کد نهایی برنامه‌ی مثال یک به صورت زیر است:

\begin{codebox}{java}
import javax.swing.*;
import java.awt.*;
import java.awt.event.*;
‌
public class Main{
   public static void main(String[] args) {  
      JFrame frame = new JFrame("Student Registration Form");  
      frame.setSize(500, 500);  
      frame.setDefaultCloseOperation(JFrame.EXIT_ON_CLOSE);  
      frame.setLayout(null);  
  
      JPanel mainPanel = new JPanel();  
      mainPanel.setBounds(40, 35, 400, 400);  
      mainPanel.setBackground(new Color(0x89CFF0));
      mainPanel.setLayout(null);  
      JPanel successPanel = new JPanel();  
      successPanel.setBounds(40, 35, 400, 400);  
      successPanel.setBackground(Color.lightGray);
      successPanel.setLayout(null);

      //labels of mainPanel  
      JLabel nameLabel = new JLabel("Student name:");  
      nameLabel.setBounds(10, 10, 150, 50);  
      JLabel numberLabel = new JLabel("Student number:");  
      numberLabel.setBounds(10, 90, 150, 50);  
      JLabel emailLabel = new JLabel("Email:");  
      emailLabel.setBounds(10, 170, 150, 50);

      //label of successPanel  
      JLabel infoLabel = new JLabel("Information Submitted successfully");  
      infoLabel.setBounds(100, 170, 350, 50);  


      JTextField nameTextField = new JTextField();  
      nameTextField.setBounds(170, 10, 200, 40);  
      JTextField numberTextField = new JTextField();  
      numberTextField.setBounds(170, 90, 200, 40);  
      JTextField emailTextField = new JTextField();  
      emailTextField.setBounds(170, 170, 200, 40);  

      ImageIcon imageIcon = new ImageIcon("src/imagee.png");  
      Image scaledImage = imageIcon.getImage().getScaledInstance(150, 150, Image.SCALE_SMOOTH);  
      JLabel imageLabel = new JLabel(new ImageIcon(scaledImage));  
      imageLabel.setBounds(30, 230, 150, 150);  
  
      JButton submitButton = new JButton("Submit"); //for mainPanel  
      submitButton.setBounds(270, 330, 100, 50);   
      JButton okButton = new JButton("OK"); //for successPanel  
      okButton.setBounds(150, 330, 100, 50);


      submitButton.addActionListener(new ActionListener() {
         @Override
         public void actionPerformed(ActionEvent e){
            frame.getContentPane().removeAll();
            frame.add(successPanel);
            frame.revalidate();
            frame.repaint();  
         }  
      });  
      okButton.addActionListener(new ActionListener() {
         @Override
         public void actionPerformed(ActionEvent e){
            frame.getContentPane().removeAll();
            frame.add(mainPanel);
            frame.revalidate();
            frame.repaint();  
         }  
      });  

      mainPanel.add(imageLabel);  
      mainPanel.add(nameLabel);  
      mainPanel.add(numberLabel);  
      mainPanel.add(emailLabel);  
      mainPanel.add(nameTextField);  
      mainPanel.add(numberTextField);  
      mainPanel.add(emailTextField);  
      mainPanel.add(submitButton);

      successPanel.add(infoLabel);

      successPanel.add(okButton);  

      frame.add(mainPanel);  
      frame.setVisible(true);  
   }  
}
\end{codebox}

\textbf{مثال دو}

از \lr{\texttt{JOptionPane}} برای نمایش پاپ آپی (\lr{\texttt{dialog}})
استفاده می‌کنیم.

به کد زیر که در آن از این کلاس استفاده کردیم، دقت کنید:

\begin{codebox}{java}
submitButton.addActionListener(new ActionListener() {
   @Override
   public void actionPerformed(ActionEvent e) {  
      JOptionPane.showMessageDialog(frame, “Information submitted successfully”);
   }
});
\end{codebox}

خروجی کد بعد از کلیک روی دکمه‌ی \lr{\texttt{Submit}}:

\begin{figure}[H]
\centering
\includegraphics[width=0.8\textwidth]{images/graphics-swing-13.png}
\end{figure}

\textbf{مثال سه}

در این مثال، از متدهای \lr{\texttt{setText}} و \lr{\texttt{getText}} نیز
استفاده خواهیم کرد:

\begin{itemize}
\tightlist
\item
  متد \lr{\texttt{setText(String text)}}: این متد برای تنظیم کردن متن
  داخل یک کامپوننت گرافیکی به کار می‌آید. یعنی وقتی بخواهید یک متن خاص را
  در \lr{\texttt{JTextField}} به عنوان متن پیش‌فرض یا راهنمای متن ورودی
  نمایش دهید، می‌توانید از این متد استفاده کنید.
\item
  متد \lr{\texttt{getText()}}: این متد نیز برای دریافت متن از داخل یک
  کامپوننت گرافیکی است. مثلا اگر بخواهید بفهمید کاربر چه چیزی وارد کرده،
  می‌توانید از این متد استفاده کنید.
\end{itemize}

حالا می‌خواهیم در یک نمونه‌ی ساده‌تر از برنامه‌ی قبلی خود، از این دو متد
استفاده کنیم، به طوری که کاربر اطلاعات خود را وارد می‌کند و با زدن دکمه
\lr{\texttt{Submit}}، اطلاعات در یک \lr{\texttt{TextField}} دیگر نمایش
داده می‌شود. برای این کار، ابتدا یک \lr{\texttt{TextField}} دیگر می‌سازیم
که بعد از فشردن دکمه \lr{\texttt{Submit}}، فقط اطلاعات را نشان دهد و
دیگر نتواند ویرایش شود. به همین دلیل، ویژگی
\lr{\texttt{setEditable(false)}} را برای آن تنظیم می‌کنیم که این فیلد فقط
برای نمایش استفاده شود. در ادامه دکمه \lr{\texttt{Submit}} را نیز تعریف
می‌کنیم. کد مربوطه به شکل زیر خواهد بود:

\begin{codebox}{java}
JButton submitButton = new JButton("Submit");  
submitButton.setBounds(150, 160, 100, 30);

JTextField outputTextField = new JTextField();  
outputTextField.setBounds(10, 210, 360, 30);  
outputTextField.setEditable(false);
\end{codebox}

چون عملیات ``دریافت ورودی'' و ``نمایش خروجی'' باید بعد از کلیک روی دکمه
\lr{\texttt{Submit}} انجام شود، این عملیات‌ها را داخل متد
\lr{\texttt{actionPerformed()}} قرار می‌دهیم. این متد فقط زمانی اجرا
می‌شود که کاربر روی دکمه کلیک کند. (کدهای مربوط به فعال کردن دکمه قبلا
توضیح داده شده) حالا به کد زیر و کامنت‌های داخل آن دقت کنید:

\begin{codebox}{java}
submitButton.addActionListener(new ActionListener() {
   @Override
   public void actionPerformed(ActionEvent e) {

      // Receive input from the name field  
      String name = nameTextField.getText();

      // Receive input from the student number field  
      String number = numberTextField.getText();

      // Receive input from the email field  
      String email = emailTextField.getText();

      // Combine the information and display it in the output field
      outputTextField.setText("Name: " + name + " | Number: " + number + " | Email: " + email);  
   }  
});
\end{codebox}

در کد بالا ما یک \lr{\texttt{ActionListener}} به دکمه اضافه کردیم که با
هر بار کلیک کردن، متد \lr{\texttt{()actionPerformed}} را اجرا کند. حالا
باید دکمه‌ی \lr{\texttt{Submit}} و \lr{\texttt{outputTextField}} را به
پنل اضافه کنیم تا همه چیز به درستی نمایش داده شود. کد مربوط به اضافه
کردن دکمه و \lr{\texttt{outputTextField}} به پنل به صورت زیر است:

\begin{codebox}{java}
mainPanel.add(submitButton);  
mainPanel.add(outputTextField);
\end{codebox}

کد نهایی به صورت زیر هست:

\begin{codebox}{java}
import javax.swing.*;
import java.awt.*;
import java.awt.event.ActionEvent;
import java.awt.event.ActionListener;

public class Main {
   public static void main(String[] args) {  
      JFrame frame = new JFrame("Student Registration Form");  
      frame.setSize(500, 500);  
      frame.setDefaultCloseOperation(JFrame.EXIT_ON_CLOSE);  
      frame.setLayout(null);  
  
      JPanel mainPanel = new JPanel();  
      mainPanel.setBounds(40, 35, 400, 400);  
      mainPanel.setBackground(new Color(0x89CFF0));  
      mainPanel.setLayout(null);  
  
      JLabel namelabel = new JLabel("Student name:");  
      namelabel.setBounds(10, 10, 150, 30);  
      JLabel numberlabel = new JLabel("Student number:");  
      numberlabel.setBounds(10, 60, 150, 30);  
      JLabel emaillabel = new JLabel("Email:");  
      emaillabel.setBounds(10, 110, 150, 30);  
  
      JTextField nametextField = new JTextField();  
      nametextField.setBounds(170, 10, 200, 30);  
      JTextField numbertextField = new JTextField();  
      numbertextField.setBounds(170, 60, 200, 30);  
      JTextField emailtextField = new JTextField();  
      emailtextField.setBounds(170, 110, 200, 30);  
  
      JButton submitButton = new JButton("Submit");  
      submitButton.setBounds(150, 160, 100, 30);  
  
      JTextField outputTextField = new JTextField();  
      outputTextField.setBounds(10, 210, 360, 30);  
      outputTextField.setEditable(false); 

      submitButton.addActionListener(new ActionListener() {
         @Override
         public void actionPerformed(ActionEvent e) {  
            String name = nametextField.getText();  
            String number = numbertextField.getText();  
            String email = emailtextField.getText();

            outputTextField.setText("Name: " + name + " | Number: " + number + " | Email: " + email);  
         }  
      });  
  
      mainPanel.add(nametextField);  
      mainPanel.add(numbertextField);  
      mainPanel.add(emailtextField);  
      mainPanel.add(namelabel);  
      mainPanel.add(numberlabel);  
      mainPanel.add(emaillabel);  
      mainPanel.add(submitButton);  
      mainPanel.add(outputTextField);  

      frame.add(mainPanel);  
      frame.setVisible(true);  
   }  
}
\end{codebox}

خروجی کد بعد از کلیک کردن بر روی دکمه:

\begin{figure}[H]
\centering
\includegraphics[width=0.5\textwidth]{images/graphics-swing-14.png}
\end{figure}

\subsection{\lr{Layouts} (چیدمان‌ها)}

چیدمان کامپوننت‌ها در صفحه، یکی از بخش‌های مهم برای ساخت رابط کاربری به
قولی \lr{user friendly} تر است. همانطور که در کد نمونه‌ی اول دیدید:

\begin{codebox}{java}
frame.setLayout(null);
\end{codebox}

این خط کد یک صفحه‌ی خالی برای ما ایجاد کرد و این امکان را داد که با مشخص
کردن موقعیت و اندازه‌ی اجزا، چیدمان صفحه را به سلیقه‌ی خودمان تنظیم کنیم.
اما یک مشکل اساسی دارد! همانطور که احتمالا خودتان هم متوجه شدید، وقتی
اندازه‌ی صفحه تغییر می‌کند، عناصر صفحه به‌صورت خودکار تنظیم نمی‌شوند و همان
جا که بودند، ثابت می‌مانند. این موضوع می‌تواند باعث به‌هم‌ریختگی رابط کاربری
شود.

برای حل این مشکل، می‌توانیم از چیدمان‌های آماده (\lr{\texttt{Layouts}})
برای \lr{\texttt{Frame}} یا سایر اجزای صفحه استفاده کنیم. در ادامه، چند
تا از این ابزارهای آماده‌ی مدیریت چیدمان را بررسی می‌کنیم.

مزیت استفاده از \lr{\texttt{layout}}‌ها این است که برخلاف روش دستی، وقتی
اندازه‌ی پنجره تغییر می‌کند، موقعیت و اندازه‌ی اجزا نیز به‌طور خودکار تنظیم
می‌شوند و تجربه‌ی کاربری بهتری رقم می‌خورد.

چیدمان‌ها تعیین می‌کنند که عناصر اضافه‌شده به \lr{\texttt{panel}} یا
\lr{\texttt{frame}} چگونه کنار هم قرار بگیرند. در اینجا ما با
\lr{\texttt{panel}} کار می‌کنیم، ولی معمولا چیدمان‌ها روی
\lr{\texttt{frame}} پیاده‌سازی می‌شوند. شما نیز بسته به نیاز و خلاقیت
خودتان، می‌توانید هر کدام از این روش‌ها را امتحان کنید.

\subsection{\lr{Border Layout}}

حالت پیش‌فرض جاوا برای چیدمان پنل‌هاست که از 5 ناحیه تشکیل شده و شکل کلی
آن به صورت زیر است:

\begin{figure}[H]
\centering
\includegraphics[width=0.5\textwidth]{images/graphics-swing-18.png}
\end{figure}

الان برای استفاده از این چیدمان، باید برخلاف کاری که در برنامه‌ای که با
هم ساختیم انجام دادیم، ورودی متد \lr{\texttt{setLayout()}} را به جای
\lr{\texttt{null}}، شی‌ای جدید از کلاس \lr{\texttt{BorderLayout}}
بگذاریم:

\begin{codebox}{java}
mainPanel.setLayout(new BorderLayout);
\end{codebox}

برای هر کدام از اجزایی که به پنل اضافه می‌کنیم، باید اشاره کنیم که در
کدام ناحیه قرار می‌گیرد که به این منظور به شکل زیر عمل می‌کنیم:

\begin{codebox}{java}
// add a new JButton with name "NORTH" and it is on top of the container JButton northButton = new JButton("NORTH"); 

mainPanel.add(northButton, BorderLayout.NORTH);
\end{codebox}

\href{https://www.youtube.com/watch?v=PD6pd6AMoOI}{این ویدئو} را برای
درک بهتر این چیدمان بررسی کنید.

\subsection{\lr{Box Layout}}

توی این چیدمان، اجزا می‌توانند به دو شکل مرتب شوند:

\begin{enumerate}
\def\labelenumi{\arabic{enumi}.}
\tightlist
\item
  \lr{\texttt{Page Axis}} ( از بالا به پایین، مثل یک ستون)
\item
  \lr{\texttt{Line Axis}} ( از چپ به راست، مثل یک سطر)
\end{enumerate}

به بیان ساده، شما این امکان را دارید که کامپوننت‌های خود را یا زیر هم یا
کنار هم قرار دهید.

هنگام ساخت یک آبجکت از کلاس \lr{\texttt{BoxLayout}} ، باید دو تا ورودی
برای آن مشخص کنید:

اولین ورودی فریم (\lr{\texttt{JFrame}}) یا پنلی (\lr{\texttt{JPanel}})
که می‌خواهید این چیدمان روی آن اعمال شود.

دومین ورودی یک مقدار ثابت است که مشخص می‌کند چیدمان به‌صورت ستونی
(\lr{\texttt{PAGE\_AXIS}}) یا سطری (\lr{\texttt{LINE\_AXIS}}) باشد.

علاوه بر این، \lr{\texttt{BoxLayout}} این امکان را به شما می‌دهد که بین
اجزا فضای خالی ایجاد کنید تا چیدمان مرتب‌تر شود.

مثال زیر نحوه‌ی استفاده از \lr{\texttt{Page Axis}} را نشان می‌دهد:

\begin{codebox}{java}
mainPanel.setLayout(new BoxLayout(mainPanel, BoxLayout.PAGE_AXIS));
\end{codebox}

\begin{figure}[H]
\centering
\includegraphics[width=0.5\textwidth]{images/graphics-swing-19.png}
\end{figure}

به عنوان مثال، اکنون پنل اول برنامه خودمان که چیدمانش \lr{\texttt{null}}
بود را با \lr{\texttt{BoxLayout}} و \lr{\texttt{PAGE\_AXIS}} پیاده‌سازی
می‌کنیم:

\begin{codebox}{java}
mainPanel.setLayout(new BoxLayout(mainPanel, BoxLayout.PAGE_AXIS));
\end{codebox}

خروجی این چیدمان برای برنامه‌ای که ساختیم به این صورت خواهد بود:

\begin{figure}[H]
\centering
\includegraphics[width=0.5\textwidth]{images/graphics-swing-20.png}
\end{figure}

دقت شود که وقتی از \lr{\texttt{Layouts}} استفاده می‌کنیم (بر خلاف وقتی که
چیدمان را \lr{\texttt{null}} می‌گذاشتیم)، ترتیب اضافه شدن ‌کامپوننت‌ها به
پنل حائز اهمیت می‌شود. مثلا برای خروجی بالا، ترتیب بدین صورت است:

\begin{codebox}{java}
mainPanel.add(nameLabel);  
mainPanel.add(nameTextField);  
mainPanel.add(numberLabel);  
mainPanel.add(numberTextField);  
mainPanel.add(emailLabel);  
mainPanel.add(emailTextField);  
mainPanel.add(imageLabel);  
mainPanel.add(submitButton);
\end{codebox}

\href{https://www.youtube.com/watch?v=hBe2eBorQuw}{این ویدئو}، برای درک
بهتر این \lr{\texttt{Layout}} پیشنهاد می‌شود.

\subsection{\lr{Flow Layout}}

چیدمان \lr{\texttt{FlowLayout}} به شما این امکان را می‌دهد تا اجزا را به
صورت پشت سر هم در یک خط یا چند خط مرتب کنید. وقتی صفحه‌تان پر شود، اجزا
به صورت خودکار به خط بعدی منتقل می‌شوند.

\begin{figure}[H]
\centering
\includegraphics[width=0.5\textwidth]{images/graphics-swing-21.png}
\end{figure}

برای درک شهودی بهتر از پیاده‌سازی این \lr{\texttt{layout}}، به
\href{https://www.youtube.com/watch?v=pDqjHozkMBs}{این ویدیو} مراجعه
کنید.

\subsection{\lr{Grid Layout}}

این نوع چیدمان به شما اجازه می‌دهد اجزا را داخل یک شبکه مرتب کنید، طوری
که هر کدام بتوانند یک یا چند سلول را اشغال کنند. شما همچنین کنترل کاملی
روی ترتیب و فضای بین این اجزا دارید.

در این روش، کامپوننت‌ها داخل یک شبکه شطرنجی قرار می‌گیرند و هر کدام، تمام
فضای سلول خودشان را پر می‌کنند. هنگام تنظیم این چیدمان، می‌توانید تعداد
سطرها، ستون‌ها و فاصله‌ی بینشان را مشخص کنید تا دقیقا همان چیزی که
می‌خواهید، پیاده‌سازی شود. همچنین می‌توانید این کارها را به ترتیب با متدهای
زیر انجام دهید:

\begin{enumerate}
\def\labelenumi{\arabic{enumi}.}
\tightlist
\item
  \lr{\texttt{setRows(int rows)}} : تنظیم تعداد سطرها
\item
  \lr{\texttt{setColumns(int coumns)}} : تنظیم تعداد ‌ستون‌ها
\item
  \lr{\texttt{setHgap(int hgap)}} : تنظیم فاصله افقی بین سلول‌ها
\item
  \lr{\texttt{setVgap(int vgap)}} : تنظیم فاصله عمودی بین سلول‌ها
\end{enumerate}

برای نمونه، در کد زیر، بر روی یک پنل با چیدمان \lr{\texttt{Grid Layout}}
شش دکمه با فاصله‌های مشخص از هم قرار می‌دهیم:

\begin{codebox}{java}
import javax.swing.*;  
import java.awt.*;  
  
public class Grid {
   public static void main(String args[]){  
      JFrame frame = new JFrame();  
      frame.setSize(500, 500);  
      frame.setDefaultCloseOperation(WindowConstants.EXIT_ON_CLOSE);  

      JPanel panel = new JPanel();
      GridLayout gridLayout = new GridLayout(2, 3, 10, 10);  
      panel.setLayout(gridLayout);  

      JButton button1 = new JButton("1");  
      JButton button2 = new JButton("2");  
      JButton button3 = new JButton("3");  
      JButton button4 = new JButton("4");  
      JButton button5 = new JButton("5");  
      JButton button6 = new JButton("6");  
  
      panel.add(button1);  
      panel.add(button2);  
      panel.add(button3);  
      panel.add(button4);  
      panel.add(button5);  
      panel.add(button6);  
  
      frame.add(panel);  
      frame.setVisible(true);  

   }
\end{codebox}

در کد بالا، یک آبجکت از کلاس \lr{\texttt{GridLayout}} ساختیم و ورودی‌ها
به ترتیب از چپ، 2 سطر، 3 ستون و فاصله‌های افقی و عمودی 10 پیکسلی بین هر
دکمه را مشخص می‌کنند.

خروجی کد:

\begin{figure}[H]
\centering
\includegraphics[width=0.5\textwidth]{images/graphics-swing-15.png}
\end{figure}

حال، با فراخوانی چهار متد زیر، تعداد سطر و ستون‌ها را برعکس کرده و
فاصله‌های عمودی و افقی را به 50 افزایش می‌دهیم:

\begin{codebox}{java}
gridLayout.setRows(3);  
gridLayout.setColumns(2);  
gridLayout.setHgap(50);  
gridLayout.setVgap(50);
\end{codebox}

خروجی کد:

\begin{figure}[H]
\centering
\includegraphics[width=0.5\textwidth]{images/graphics-swing-16.png}
\end{figure}

برای مثال، در پیاده‌سازی این برنامه نیز از \lr{grid layout} استفاده شده:

\begin{figure}[H]
\centering
\includegraphics[width=0.5\textwidth]{images/graphics-swing-17.png}
\end{figure}

\href{https://www.youtube.com/watch?v=ohNqQagkDDY}{این ویدئو} را برای
بررسی دقیق‌تر این \lr{\texttt{Layout}} نگاه کنید.

\section{برخی از متدهای پرکاربرد برای کلاس‌های کاربردی \lr{Swing}}

\subsection{کلاس \texttt{JFrame}}

\begin{itemize}
\item
  \lr{\texttt{setText(String title)}} : این متد، همان تایتل فریم را
  تغییر میدهد.
\item
  \lr{\texttt{setSize(int width, int height)}} : ابعاد فریم را مشخص
  می‌کند.
\item
  \lr{\texttt{setVisible(boolean state)}} : وضعیت نمایش فریم را مشخص
  می‌کند.
\item
  \lr{\texttt{setResizable(boolean state)}} : اگر ورودی را \lr{false}
  قرار دهیم، ابعاد فریم ثابت باقی مانده و کاربر نمی‌تواند آن را تغییر
  دهد.
\item
  \lr{\texttt{setLayout()}} : چیدمان فریم را مشخص می‌کند.
\item
  \lr{\texttt{setDefaultCloseOperation(int operation)}} : بسته به نوع
  ورودی، وضعیت برنامه بعد از بستن فریم را مشخص می‌کند. شما می‌توانید
  \lr{\texttt{operation}} را به جای \lr{\texttt{int}}، با
  \lr{\texttt{constant}}‌های از پیش تعریف شده‌ای مثل
  \lr{\texttt{JFrame.EXIT\_ON\_CLOSE}} پر کنید.
\item
  \lr{\texttt{setLocationRelativeTo(Component c)}} : با
  \lr{\texttt{null}} قرار دادن ورودی این متد، هنگام ران کردن برنامه،
  فریم در وسط صفحه‌ی مانیتور باز می‌شود.
\end{itemize}

\subsection{کلاس \lr{JLabel}}

\begin{itemize}
\item
  \lr{\texttt{setText(String text)}} : متن لیبل را مشخص می‌کند.
\item
  \lr{\texttt{setIcon()}} : مثلا برای اضافه کردن \lr{\texttt{imageIcon}}
  همانطور که بالاتر گفتیم.
\item
  \lr{\texttt{setForeGround()}} : تعیین رنگ متن لیبل به کمک کلاس
  \lr{\texttt{Color}} که بالاتر استفاده ازش را توضیح دادیم.
\item
  \lr{\texttt{setFont()}} : فونت متن لیبل را مشخص می‌کند. ورودی آن یک
  آبجکت از کلاس فونت است که به صورت زیر ساخته می‌شود:
\end{itemize}

\begin{codebox}{java}
Font font = new Font("font's name", Font.type, int size);
\end{codebox}

قسمت \lr{\texttt{type}} میتواند ثابت‌های مختلفی مثل \lr{\texttt{PLAIN}} و
یا \lr{\texttt{BOLD}} باشد.

\begin{itemize}
\tightlist
\item
  \lr{\texttt{setBorder()}}
\end{itemize}

یک مرز و قاب به دور برچسب ایجاد می‌کند. ورودی آن یک آبجکت از کلاس
\lr{\texttt{Border}} است.

برای استفاده از متد \lr{\texttt{setBorder}} مثل زیر عمل می‌کنیم:

ورودی دوم ضخامت مرز را مشخص می‌کند.

\begin{codebox}{java}
Border border = BorderFactory.createLineBorder(Color.RED, 10);  
namelabel.setBorder(border);
\end{codebox}

\begin{figure}[H]
\centering
\includegraphics[width=0.5\textwidth]{images/graphics-swing-22.png}
\end{figure}

\subsection{کلاس \lr{JButton}}

\begin{itemize}
\item
  \lr{\texttt{setText(String text)}} : مثل \lr{\texttt{JLable}} کار
  می‌کند.
\item
  \lr{\texttt{setBackGround()}} : مثل متد \lr{\texttt{setForeGround()}}
  که برای کلاس \lr{\texttt{JLable}} بود، این متد هم رنگ پس زمینه‌ی دکمه
  را با اسفاده از کلاس \lr{\texttt{Color}} مشخص می‌کند.
\item
  \lr{\texttt{setFocusable(boolean state)}} : باعث فعالسازی و یا عدم
  فعالسازی مرز کمرنگ دور نوشته‌ی درون دکمه می‌شود.
\item
  \lr{\texttt{setBorder()}} : مثل \lr{\texttt{JLabel}} کار می‌کند.
\item
  \lr{\texttt{setActionListener()}} : توی مثال‌ها توضیح داده شده است.
\end{itemize}
	
	
	
	
%%%%%%%%%%%%%%%%%%%%%%%%%%%%%%%%%%%%%%%%%%%%%%%%%%%%%%%%%%%%
\end{document}